% paper_geovac_lcao_fci.tex — GeoVac LCAO FCI Paper
% Topological Full Configuration Interaction for Heteronuclear Diatomics
% Status: Scaffold, March 2026
% Target: arXiv (physics.comp-ph / quant-ph)
\documentclass[aps,pra,twocolumn,superscriptaddress,longbibliography]{revtex4-2}

\usepackage{amsmath,amssymb,amsthm,graphicx,xcolor,bm,braket,dcolumn,booktabs}

\newtheorem{theorem}{Theorem}
\newtheorem{corollary}{Corollary}

\begin{document}

\title{Topological Full Configuration Interaction for Heteronuclear\\
Diatomics: LiH as a Benchmark}
\thanks{GeoVac FCI Paper in the Geometric Vacuum series}

\author{J.~Loutey}
\affiliation{Independent Researcher}

\date{\today}

\begin{abstract}
We present a topological full configuration interaction (FCI) method for
heteronuclear diatomic molecules based on the GeoVac discrete graph
Laplacian framework.  Two atom-centered hydrogenic lattices on $S^3$ are
coupled via Mulliken cross-nuclear attraction, Slater--Condon electron
repulsion integrals, and STO-overlap bridges to form a molecular
Hamiltonian diagonalized in the full Slater determinant space.  For LiH
at nmax$= 3$ (28 spatial orbitals, 56 spin-orbitals, 367{,}290 Slater
determinants), the Boys--Bernardi counterpoise-corrected dissociation
energy is $D_e^{\text{CP}} = 0.093$~Ha (1.0\% error vs.\ the
experimental $0.092$~Ha).  The equilibrium bond length
$R_{\text{eq}} \approx 2.5$~bohr (experiment: 3.015~bohr) is identified
as a discretization artifact of the nmax$= 3$ basis truncation,
predicted to converge toward experiment with increasing nmax by analogy
with the Paper~1 $s/p$ splitting precedent.

A systematic diagnostic arc of 29~versions (v0.9.8--v0.9.36) tested
and eliminated every specific hypothesis for the $R_{\text{eq}}$ error:
exact cross-nuclear attraction (shell theorem), $l > 0$ cross-atom
electron repulsion ($+0.449$~mHa, negligible), orbital exponent
relaxation ($R$-independent shift), and the full Sturmian bond sphere
basis (structural theorem: Hamiltonian proportional to overlap, eigenvalues
geometry-independent).  The $SO(4)$ Wigner $D$-matrix selection rules
$D^2_{10,10}(\gamma) = 1$ and $D^2_{00,10}(\gamma) = 0$ are proven as
structural results of the bond sphere geometry.

The excitation-driven direct CI solver achieves $O(N_{\text{SD}} \times
N_{\text{connected}})$ scaling, enabling four-electron FCI calculations
(Be nmax$= 4$: 487{,}635 SDs in 357~s, 0.90\% error).
\end{abstract}

\maketitle


%% ========================================================================
\section{Introduction}
\label{sec:intro}
%% ========================================================================

% PLACEHOLDER — to be written next session
%
% Motivation: quantum chemistry from discrete graph topology.
% GeoVac premise from Paper 0: quantum numbers (n,l,m) are geometric
% coordinates on S^3. The graph Laplacian is mathematically equivalent
% to the Schrodinger equation via Fock's 1935 projection (Paper 7).
%
% Why LiH: simplest heteronuclear test case, well-characterized
% experimentally (D_e = 0.092 Ha, R_eq = 3.015 bohr), four electrons
% accessible to FCI, large charge asymmetry (Z_A/Z_B = 3) tests
% the framework at its limits.
%
% Previous GeoVac results: single-electron atoms < 0.1% error (Paper 1),
% He FCI 0.35% (v0.9.5), Li FCI 1.1% (v0.9.7), H2 FCI < 2% (v0.9.4).
% This paper extends to heteronuclear molecules.
%
% Key references: Fock 1935, Papers 0, 1, 7, 8.


%% ========================================================================
\section{Architecture}
\label{sec:architecture}
%% ========================================================================

% PLACEHOLDER — to be written next session
%
% 2.1 MolecularLatticeIndex
%   Two atomic subgraphs (GeometricLattice, one per atom).
%   Combined basis: N_A + N_B spatial orbitals, 2(N_A + N_B) spin-orbitals.
%   Slater determinant space: C(2N_sp, N_e) determinants.
%   Duck-type compatible with DirectCISolver.
%
% 2.2 One-Electron Hamiltonian
%   H1 = H_A tensor I + I tensor H_B + V_cross
%   H_A, H_B: exact atomic eigenvalues on diagonal (-Z^2/(2n^2))
%   V_cross: Mulliken cross-nuclear (diagonal) + graph bridges (off-diagonal)
%   Eq: H1_total = diag(-Z_A^2/(2n^2)) + Mulliken_cross_B + bridges
%
% 2.3 Two-Electron Repulsion
%   V_ee: same-atom Slater F^0 (exact) + cross-atom Coulomb J (monopole)
%   + cross-atom exchange K (s-s only, Mulliken)
%   Eq: V_ee = sum over pairs of F^0 or J integrals
%
% 2.4 Nuclear Repulsion
%   V_NN = Z_A * Z_B / R (exact, constant shift)
%
% 2.5 Bridge Construction
%   STO overlap x conformal factors, priority: boundary states first
%   connectivity = [(atom_i, atom_j, n_bridges)]
%
% 2.6 BSSE and Counterpoise Correction
%   Boys-Bernardi 1970: E_BSSE = E(A in AB basis) - E(A in A basis)
%   Ghost atoms (Z=0): zero h1_diag, no bridges, no cross-nuclear
%   compute_bsse_correction() standalone function
%
% 2.7 O(V) Scaling from DirectCISolver
%   Excitation-driven sparse Hamiltonian: O(N_SD x N_connected)
%   Knowles-Handy 1984 algorithm
%   Dense lookup tables for hot loops (H1, ERI)


%% ========================================================================
\section{LiH Benchmark Results}
\label{sec:results}
%% ========================================================================

% PLACEHOLDER — to be written next session
%
% 3.1 CP-Corrected PES at nmax=3
%   Table: R, E_mol, E_atoms, D_e_raw, BSSE, D_e_CP
%   Key points: R = 2.0, 2.5, 3.015, 3.5, 4.0, 5.0, 6.0, 8.0, 10.0
%
% 3.2 Best Results
%   D_e_CP = 0.093 Ha at R=3.015 (1.0% error vs expt 0.092)
%   R_eq ~ 2.5 bohr (expt 3.015)
%   BSSE = -0.115 Ha (Li: -0.105, H: -0.010)
%   E_mol(R=3.015) = -7.981 Ha (raw), -7.866 Ha (CP-corrected)
%
% 3.3 Error Analysis
%   D_e accurate (1% error) because BSSE is R-independent
%   R_eq shifted because cross-nuclear attraction overbinds at short R
%   Mulliken/Fourier diagonal: V_cross_B screening deficit = -0.247 Ha
%
% 3.4 Comparison to Experiment
%   Huber-Herzberg: D_e = 0.092 Ha, R_eq = 3.015 bohr


%% ========================================================================
\section{Systematic Diagnostic Arc}
\label{sec:diagnostic}
%% ========================================================================

% PLACEHOLDER — to be written next session
%
% This is the scientifically valuable section. 29 versions, each a
% controlled experiment. Organized by hypothesis tested.

\subsection{BSSE Diagnosis and Counterpoise Correction (v0.9.8--v0.9.11)}
% v0.9.8: first LiH FCI, E = -8.097 Ha, variational violation
% v0.9.9: BSSE identified, Boys-Bernardi CP correction implemented
% v0.9.10: ghost atom infrastructure (Z=0)
% v0.9.11: CP-corrected PES, D_e_CP = 0.093 Ha (1.0% error)

\subsection{Cross-Nuclear and V\_ee Completeness (v0.9.12--v0.9.13, v0.9.35--v0.9.36)}
% v0.9.12: Lowdin 4-index ERI transform — catastrophically wrong
% v0.9.13: Mulliken exchange K — D_e_CP = 0.110 Ha (18.8% error)
% v0.9.35: l>0 cross-atom V_ee — dE = +0.449 mHa (negligible)
% v0.9.36: orbital exponent relaxation — R-independent shift

\subsection{SO(4) D-Matrix Coupling (v0.9.15--v0.9.18)}
% v0.9.15: SO(4) Wigner D-matrix implementation (44/44 tests)
% v0.9.16: geometric-mean coupling — LiH unbound
% v0.9.17: SW off-diagonal only — missing diagonal cross-nuclear
% v0.9.18: hybrid (Fourier diag + SW off-diag) — LiH bound, D_e_CP = 0.143

\subsection{Sturmian Basis Arc (v0.9.20--v0.9.34)}
% Summary from Paper 8: structural theorem proven
% All variants failed: atom-centered, molecular, MO-FCI, dual-p0
% See Paper 8 Sec XII for the structural theorem

\subsection{R\_eq Root Cause Elimination}
% v0.9.24: term decomposition — V_cross_B drives R_eq inward
% v0.9.35: l>0 V_ee — 0.449 mHa, negligible
% v0.9.36: orbital relaxation — R-independent shift
% Harmonic phase locking — no D-matrix critical points at R_eq
% Conclusion: R_eq error is nmax discretization, not missing physics


%% ========================================================================
\section{$R_{\text{eq}}$ as Discretization Artifact}
\label{sec:req}
%% ========================================================================

% PLACEHOLDER — to be written next session
%
% 5.1 The Paper 1 Analogy
%   In Paper 1, the s/p orbital splitting was initially too large at
%   small nmax and converged toward the exact value as nmax increased.
%   The mechanism: finite graph truncation creates artificial energy
%   penalties for higher angular momentum states that decrease as the
%   graph becomes denser.
%
% 5.2 The nmax Convergence Argument
%   At nmax=3, the basis has 14 spatial orbitals per atom.
%   The cross-nuclear Mulliken overlap integral uses these truncated
%   atomic wavefunctions. At short R, the overlap samples high-momentum
%   components that are missing from the nmax=3 basis, creating an
%   artificial attraction that decays too slowly with R.
%   At nmax -> infinity, the basis is complete and the overlap converges
%   to the exact Mulliken integral.
%
% 5.3 The Prediction
%   R_eq -> 3.015 bohr as nmax -> infinity.
%   The convergence rate is unknown but expected to be algebraic (not
%   exponential) in nmax, based on the Paper 1 precedent.
%
% 5.4 Cite the 2s/2p Splitting Precedent
%   Paper 1: 2s/2p splitting at nmax=3 was X% from exact, converged
%   to Y% at nmax=10. Same mechanism: basis truncation creates
%   artificial energy penalties that decrease with nmax.
%
% 5.5 What Would Confirm This
%   nmax=4 LiH FCI: if R_eq shifts toward 3.015, the prediction is
%   supported. Current limitation: nmax=4 requires ~500k SDs for LiH,
%   feasible with DirectCISolver but not yet attempted.


%% ========================================================================
\section{Conclusion}
\label{sec:conclusion}
%% ========================================================================

% PLACEHOLDER — to be written next session
%
% What works:
% - Topological LCAO FCI for heteronuclear diatomics
% - D_e_CP = 0.093 Ha (1.0% error) for LiH at nmax=3
% - Excitation-driven direct CI with O(N_SD x N_connected) scaling
% - BSSE diagnosis and counterpoise correction
%
% What the limits are:
% - R_eq shifted (discretization artifact, predicted to converge)
% - nmax=3 basis truncation limits accuracy
% - Cross-nuclear attraction uses Mulliken approximation
%
% Future work:
% - nmax=4 LiH verification (would confirm R_eq convergence prediction)
% - Larger heteronuclear systems (HF, H2O)
% - Adaptive bridge construction


%% ========================================================================
\appendix
\section{Version History Table}
\label{app:versions}
%% ========================================================================

% PLACEHOLDER — to be filled with data from the 29 versions
%
% Table format:
% Version | Method | E_mol (Ha) | D_e_CP (Ha) | R_eq (bohr) | Bound?
%
% v0.9.8  | LCAO baseline          | -8.097 | 0.198 (raw) | ~2.0 | Yes
% v0.9.9  | + counterpoise          | ...    | ...         | ...  | ...
% v0.9.11 | CP-corrected PES        | ...    | 0.093       | ~2.5 | Yes
% v0.9.12 | + Lowdin ERI            | ...    | WRONG       | ...  | ...
% v0.9.13 | + Mulliken K            | ...    | 0.110       | ~2.5 | Yes
% v0.9.15 | D-matrix impl           | ...    | ...         | ...  | ...
% v0.9.16 | D-matrix geom-mean      | ...    | ...         | ...  | No
% v0.9.17 | D-matrix SW off-diag    | ...    | ...         | ...  | No
% v0.9.18 | D-matrix hybrid         | ...    | 0.143       | <2.0 | Yes
% v0.9.20 | Sturmian basic          | ...    | ...         | ...  | No
% v0.9.21 | + SW off-diagonal       | ...    | ...         | ...  | No
% v0.9.22-25 | Various corrections  | ...    | ...         | ...  | Mixed
% v0.9.26-27 | SO(4) CG V_ee        | ...    | ...         | ...  | No
% v0.9.28 | Mol beta node weights   | ...    | ...         | ...  | No
% v0.9.29 | Prolate spheroidal      | ...    | ...         | ...  | N/A
% v0.9.30 | MO projection           | ...    | ...         | ...  | Unphys
% v0.9.31 | MO Sturmian FCI         | -10.07 | ...         | <2.0 | Over
% v0.9.32 | + exact J               | -7.131 | ...         | ~3.5 | No
% v0.9.33 | + exact K               | -6.796 | ...         | ...  | No
% v0.9.34 | Dual-p0                 | -10.1  | ...         | <2.0 | Over
% v0.9.35 | l>0 cross-atom V_ee     | ...    | 0.109       | ~2.5 | Yes
% v0.9.36 | Orbital relaxation      | ...    | 0.433       | ~2.5 | Yes*


%% ========================================================================
\begin{acknowledgments}
Computational resources were provided by personal hardware.
\end{acknowledgments}
%% ========================================================================

\begin{thebibliography}{25}

\bibitem{loutey_paper0}
J.~Loutey,
``The Geometric Atom: Quantum State Space as a Packing Problem,''
GeoVac Paper~0 (2026).

\bibitem{loutey_paper1}
J.~Loutey,
``The Geometric Atom: Quantum Mechanics as a Packing Problem,''
GeoVac Paper~1 (2026).

\bibitem{loutey_paper7}
J.~Loutey,
``The Dimensionless Vacuum: Recovering the Schr\"odinger Equation
from Scale-Invariant Graph Topology,''
GeoVac Paper~7 (2026).

\bibitem{loutey_paper8}
J.~Loutey,
``Bond Sphere: Diatomic Molecules as $S^3$ and their Natural
Sturmian Basis,''
GeoVac Paper~8--9 (2026).

\bibitem{Fock1935}
V.~Fock,
``Zur Theorie des Wasserstoffatoms,''
Z. Phys. \textbf{98}, 145--154 (1935).

\bibitem{Boys1970}
S.~F. Boys and F.~Bernardi,
``The calculation of small molecular interactions by the differences
of separate total energies. Some procedures with reduced errors,''
Mol. Phys. \textbf{19}, 553--566 (1970).

\bibitem{shull_lowdin}
H.~Shull and P.-O.~L\"owdin,
``Superposition of configurations and natural spin orbitals.
Applications to the He problem,''
J. Chem. Phys. \textbf{25}, 1035 (1956).

\bibitem{Huber1979}
K.~P. Huber and G.~Herzberg,
\textit{Molecular Spectra and Molecular Structure. IV. Constants of
Diatomic Molecules}
(Van Nostrand Reinhold, New York, 1979).

\bibitem{Knowles1984}
P.~J. Knowles and N.~C. Handy,
``A new determinant-based full configuration interaction method,''
Chem. Phys. Lett. \textbf{111}, 315--321 (1984).

\end{thebibliography}

\end{document}
